\section{Дифтонги}

Два следующих друг за другом гласных могут сливаться в один сложный, неделимый звук --- дифтонг.

Для правильного произношения дифтонгов рекомендуется произнести соответствующие гласные
сначала раздельно, но без паузы, а затем максимально слитно и плавно.

\begin{longtblr}{
    colspec = { X[1,c,t] X[1,c,t] X[1,c,t] X[4,l,t] },
    rowhead=1,
    row{1} = {font=\bfseries},
}
\toprule
Румынский звук & Русский аналог & Написание & Комментарий \\
\midrule
\textipa{[\t*{ea}]}
& --
& ea
&
\\
\textipa{[\t*{oa}]}
& --
& oa
& Встречается только в ударном положении.
\\
\bottomrule
\end{longtblr}

Примеры:
\begin{description}
    \item[\textipa{[\t*{ea}]}:] mea \textipa{[m\t*{ea}]}, seamă \textipa{["s\t*{ea}m@]};
        marea \textipa{["mar\t*{ea}]}, unirea \textipa{[u"nIr\t*{ea}]}.
    \item[\textipa{[\t*{oa}]}:] poartă \textipa{["p\t*{oa}rt@]}, noapte \textipa{["n\t*{oa}pte]},
        frumoasă \textipa{[fru"m\t*{oa}s@]}; foamea \textipa{["f\t*{oa}m\t*{ea}]},
        sudoarea \textipa{[su"d\t*{oa}r\t*{ea}]}.
\end{description}